Une enfant de \qty{\v{m}}{kg} part du repos au sommet d'une glissade de \qty{\v{h}}{m} de hauteur. La glissade mesure \qty{\v{L}}{m} de long. Au bas, elle arrive avec une vitesse de seulement \qty{\v{v}}{m/s}, à cause du frottement. On prend le bas de la glissade comme niveau de référence et $g = \qty{9.81}{m/s^2}$.

\begin{enumerate}[label=\alph*)]
    \item Quelle est l'énergie mécanique de l'enfant au sommet ?
    \item Quelle est son énergie cinétique au bas de la glissade ?
    \item Quelle quantité d'énergie a été transformée en énergie thermique par le frottement ?
    \item Quelle est la force de frottement moyenne exercée sur l'enfant le long de la glissade ?
\end{enumerate}
